\section{Changing trajectories, likelihoods, and conditional consumers}
\label{sec:v18-transport}
\subsection{A common preparation, not a common flow}
Let $(\mathcal Z,\nu)$ be a standard Borel probability space, and let
$H,H_n:[0,T]\times\mathcal Z\to\R$ be jointly measurable, with
$|H|,|H_n|\le M$. They may be observations of different flows on different
phase spaces, pulled back by specified preparation maps from $\mathcal Z$.
No common microscopic generator or invariant measure is assumed. Set
\[
 \eta_n^2=\int_{\mathcal Z}\int_0^T|H_n(t,z)-H(t,z)|^2\,dt\,d\nu(z).
\]
On $\mathcal Z\times C([0,T])$ let $R$ be the law of independent $Z\sim\nu$
and $Y=\sigma B$, where $\sigma>0$. Define
\begin{align}
 L_t^n&=\exp\left\{\frac1\sigma\int_0^t H_n(s,Z)\,dB_s
                   -\frac1{2\sigma^2}\int_0^t H_n(s,Z)^2\,ds\right\},
                  \label{eq:v18-likelihood}\\
 dP_n&=L_T^n\,dR.
\end{align}
The reference model $P,L$ uses $H$. Conditional on $Z$, these are Gaussian
shifts, so under $P_n$ the observed path has law
$Y_t=\int_0^tH_n(s,Z)ds+\sigma B_t^n$. Boundedness gives true strictly
positive density martingales. Write $\mathcal F_t$ for the usual canonical
observation filtration and $\mathcal H_t=\sigma(Z,B_s:s\le t)$ for the
larger reference filtration. Equivalent laws have the same completed
observation stopping times.

\begin{theorem}[Likelihood and posterior transport for changing paths]
\label{thm:v18-gaussian}
The measures $P_n,P,R$ are equivalent and have common $Z$-marginal $\nu$.
With $A=M^2T/\sigma^2$ and
$\varepsilon_n=\min\{1,\eta_n/(2\sigma)\}$, one has
\begin{align}
 \|P_n-P\|_{\rm TV}&\le\varepsilon_n,
                     \label{eq:v18-gaussiantv}\\
 \E_R\sup_{t\le T}|L_t^n-L_t|^2
   &\le\frac8{\sigma^2}e^{3A}\eta_n^2.
                     \label{eq:v18-densityS2}
\end{align}
For $|f(Z)|\le K$ put
\[
 N_t^{n,f}=\E_R[L_T^nf(Z)\mid\mathcal F_t],\qquad
 \ell_t^n=N_t^{n,1}.
\]
Then
\begin{equation}\label{eq:v18-projectedS2}
 \E_R\sup_t|N_t^{n,f}-N_t^f|^2
              \le\frac{8K^2}{\sigma^2}e^{3A}\eta_n^2.
\end{equation}
In particular this applies to the projected likelihood $\ell^n$.
For every $|f|\le1$, the posterior functions satisfy the same-path bound
$80\varepsilon_n$ of Theorem~\ref{thm:v17-optional}, and there is one
latent-preserving coupling, independent of $f$, on which
\begin{equation}\label{eq:v18-posteriorS2}
 \E\sup_t|\E_{P_n}[f(Z)\mid\mathcal F_t^{Y^n}]
                 -\E_P[f(Z)\mid\mathcal F_t^Y]|^2
                    \le84\varepsilon_n.
\end{equation}
If $\eta_n\to0$ and deterministic microscopic paths $X_n(Z)\to X(Z)$
in a specified Polish path metric in $\nu$-probability, then on this
coupling the microscopic paths, the full and projected likelihood
processes evaluated on their own observation paths, and each bounded
posterior process converge jointly in probability in that metric and
uniformly in time for the last three coordinates.
\end{theorem}
\begin{proof}
Conditional Gaussian change of measure gives
\[
 \KL(P_n\|P)=\frac1{2\sigma^2}\E_\nu\int_0^T|H_n-H|^2dt
                         =\frac{\eta_n^2}{2\sigma^2}.
\]
The same calculation proves conditional equivalence, and the conditional
density has mean one, so the preparation marginal is unchanged. Integrating
and applying Pinsker proves \eqref{eq:v18-gaussiantv}.

For the quantitative density estimate, condition first on $Z=z$ and write
$D_t=L_t^n-L_t$. The stochastic differential equation is
\[
 dD_t=\sigma^{-1}\{D_tH_n(t,z)+L_t(H_n-H)(t,z)\}\,dB_t.
\]
The conditional second moment of $L_t$ is at most $e^A$. It\^o's isometry
and $(u+v)^2\le2u^2+2v^2$ therefore imply
\[
 \E_R[D_t^2\mid z]
 \le\frac{2M^2}{\sigma^2}\int_0^t\E_R[D_s^2\mid z]ds
     +\frac{2e^A}{\sigma^2}\int_0^t|H_n-H|^2(s,z)ds.
\]
Gronwall, followed by integration in $z$, gives
$\E_R D_T^2\le(2/\sigma^2)e^{3A}\eta_n^2$.
The difference is a square-integrable $\mathcal H_t$-martingale.
Doob's $L^2$ inequality proves \eqref{eq:v18-densityS2}. Applying the same
maximal inequality to the $\mathcal F_t$-martingale
$\E_R[(L_T^n-L_T)f\mid\mathcal F_t]$ proves
\eqref{eq:v18-projectedS2}. These are direct process estimates, not a
claim that finite-dimensional convergence alone implies tightness.

The posterior assertions follow from Theorem~\ref{thm:v17-optional} and
its countable conditional-coupling construction. For the last claim,
on the event $Y^n=Y$ the two likelihood functions are evaluated on the
same $(Z,Y)$. Their uniform difference tends to zero in $R$-probability
by the estimates just proved, and therefore in $P$-probability by absolute
continuity of the fixed $P$ with respect to $R$. The complementary event
has probability at most $\varepsilon_n$. The same argument applies to
$\ell^n$, and \eqref{eq:v18-posteriorS2} handles the posterior. The
microscopic convergence holds on any coupling retaining the same $Z$.
A finite union bound proves joint convergence. Stochastic integral versions
are chosen measurably on the canonical product space; equivalence makes
these common versions valid under every law in the countable sequence.
\end{proof}

For $f\in L^2(\nu)$ the posterior conclusion follows by truncation exactly
as in Corollary~\ref{cor:v17-changing}; it also permits $f_n\to f$ in
$L^2(\nu)$. Boundedness of $H_n$ is being used for the uniform likelihood
second moments. It is not replaced by an unverified uniform-integrability
assertion in a singular or zero-noise limit.

\subsection{The projected stochastic exponential}
Let $\widehat H_t^n=\E_{P_n}[H_n(t,Z)\mid\mathcal F_t]$, using a predictable
version for Lebesgue-almost every $t$. Put $\theta_t^n=\widehat H_t^n/\sigma$.

\begin{proposition}[Identification before passage to the limit]
\label{prop:v18-innovation}
The projected likelihood and the innovation satisfy
\begin{align}
 \ell_t^n&=\exp\left\{\int_0^t\theta_s^n\,dB_s
                         -\tfrac12\int_0^t(\theta_s^n)^2ds\right\},
                         \label{eq:v18-projectedexp}\\
 I_t^n&=B_t-\int_0^t\theta_s^n ds.
\end{align}
Here $I^n$ is a $P_n$-Brownian motion in $\mathcal F_t$. In particular the limit of the
projected likelihoods in \eqref{eq:v18-projectedS2} is the projected
stochastic exponential of the limiting model. Change of measure commutes
with this limit on bounded measurable path tests.
\end{proposition}
\begin{proof}
Project $L_t^n=1+\int_0^t L_s^nH_n(s,Z)\,dB_s/\sigma$ onto the Brownian
observation filtration. The integrands are square integrable; testing
against bounded elementary predictable stochastic integrals and using
It\^o's isometry gives
\[
 d\ell_t^n=\sigma^{-1}\E_R[L_t^nH_n(t,Z)\mid\mathcal F_t]dB_t
                 =\ell_t^n\theta_t^n dB_t.
\]
The last equality is Bayes' formula. Since $|\theta_t^n|\le M/\sigma$,
the solution is the positive stochastic exponential
\eqref{eq:v18-projectedexp}; Girsanov gives the stated innovation. The
same argument with $H$ identifies the limiting exponential. Finally
$|\E_{P_n}F-\E_PF|\le2\|F\|_\infty\varepsilon_n$ for bounded signed
$F$. There is no interchange based only on a formal limit of logarithms.
\end{proof}

\subsection{A bounded entropic backward-equation consumer}
Fix $\gamma>0$ and a latent payoff $f$ with $|f|\le1$. Define
\begin{equation}\label{eq:v18-entropicvalue}
 U_t^n=\frac1\gamma\log\E_{P_n}[e^{\gamma f(Z)}\mid\mathcal F_t]
       =\frac1\gamma\log\frac{N_t^{n,e^{\gamma f}}}{\ell_t^n}.
\end{equation}
The terminal condition here is the observable conditional certainty
equivalent $U_T^n$, not the generally unobserved $f(Z)$ itself. This
specification avoids imposing a latent terminal value on the smaller
observation filtration.

\begin{theorem}[Entropic values and localized backward integrands]
\label{thm:v18-entropic}
There is a unique bounded entropic value process with terminal value
$U_T^n$; for a predictable integrand $Z^n$ it satisfies
\begin{equation}\label{eq:v18-bsde}
 U_t^n=U_T^n+\int_t^T\frac\gamma2|Z_s^n|^2ds
                        -\int_t^T Z_s^n\,dI_s^n.
\end{equation}
On the coupling of Theorem~\ref{thm:v18-gaussian},
\begin{equation}\label{eq:v18-entropiccoupling}
 \E\sup_t|U_t^n(Y^n)-U_t(Y)|^2
       \le84\left(\frac{e^{2\gamma}-1}{\gamma}\right)^2
                                                     \varepsilon_n.
\end{equation}
For each $r>e^\gamma$ let $\tau_r^n$ be the first exit of any of
$\ell^n,\ell,N^{n,e^{\gamma f}},N^{e^{\gamma f}}$ from $(1/r,r)$, capped
at $T$. If $\eta_n\to0$, then
\begin{equation}\label{eq:v18-bsdeintegrands}
 \E_R\int_0^{\tau_r^n}|Z_s^n-Z_s|^2ds\longrightarrow0,
 \qquad
 \E_R\int_0^{\tau_r^n}|\theta_s^n-\theta_s|^2ds\longrightarrow0.
\end{equation}
Furthermore $\lim_{r\to\infty}\sup_n R(\tau_r^n<T)=0$. Thus the
localizations exhaust the observation experiment in probability.
\end{theorem}
\begin{proof}
For a positive bounded latent $g$, write
$N_t^{n,g}=\E_R[L_T^ng\mid\mathcal F_t]$ and
$dN_t^{n,g}=a_t^{n,g}dB_t$. Projection of the full integral gives
\[
 a_t^{n,g}=\sigma^{-1}\E_R[L_t^ngH_n(t,Z)\mid\mathcal F_t],
 \qquad |a_t^{n,g}|\le(M/\sigma)N_t^{n,g}.
\]
The same representation with $g=1$ gives $d\ell_t^n=b_t^n dB_t$ and
$b_t^n/\ell_t^n=\theta_t^n$. Put $N^n=N^{n,e^{\gamma f}}$ and
$a^n=a^{n,e^{\gamma f}}$. It\^o's formula in
\eqref{eq:v18-entropicvalue} gives the martingale coefficient
\[
 Z_t^n=\gamma^{-1}\left(a_t^n/N_t^n-b_t^n/\ell_t^n\right).
\]
Replacing $dB_t$ by $dI_t^n+\theta_t^n dt$ reduces its drift to
$-\gamma|Z_t^n|^2/2$, which proves \eqref{eq:v18-bsde}. The ratios
are bounded by $M/\sigma$, so the integrals are well defined. Conversely,
exponentiating any bounded solution of \eqref{eq:v18-bsde} produces a
bounded $P_n$-martingale with terminal value $e^{\gamma U_T^n}$.
Its conditional-expectation representation is unique, and martingale
representation then identifies its integrand. This proves uniqueness in
the bounded class.

Apply Theorem~\ref{thm:v17-optional} to
$g_0=(e^{\gamma f}-e^{-\gamma})/(e^\gamma-e^{-\gamma})\in[0,1]$.
The map taking its conditional expectation to
$\gamma^{-1}\log\E[e^{\gamma f}\mid\mathcal F_t]$ has Lipschitz constant
at most $(e^{2\gamma}-1)/\gamma$. This proves
\eqref{eq:v18-entropiccoupling}.

The initial values of $N^n$ and $N$ agree, since $Z$ has the same
marginal. It\^o's isometry and the terminal version of
\eqref{eq:v18-projectedS2} imply
\[
 \E_R\int_0^T|a_s^n-a_s|^2ds\to0,
 \qquad \E_R\int_0^T|b_s^n-b_s|^2ds\to0.
\]
Before $\tau_r^n$, all denominators lie in $[1/r,r]$, and the limiting
numerators $a,b$ are bounded in absolute value by $(M/\sigma)r$.
For example,
\[
 \left|\frac{a^n}{N^n}-\frac aN\right|
 \le r|a^n-a|+(M/\sigma)r^3|N^n-N|.
\]
The supremum estimates \eqref{eq:v18-projectedS2}, followed by integration
in time, prove \eqref{eq:v18-bsdeintegrands}. The same inequality applies
to $b^n/\ell^n$.

Finally $e^{-\gamma}\ell^n\le N^n\le e^\gamma\ell^n$. Each $\ell^n$
is a stochastic exponential with uniformly bounded integrand
$\theta^n$. Its positive and negative moments of the supremum are
bounded uniformly in $n$: write any power as a stochastic exponential
times its bounded compensator and apply Doob's inequality to a higher
positive power. The same calculation applies to its inverse. Markov's
inequality, the preceding comparison of $N^n$ with $\ell^n$, and a union
bound prove the exhaustion assertion. This is a localized integrand
theorem; it does not assert an unspecified global topology for backward
equations with unbounded terminal data or different nonlinear drivers.
\end{proof}

\begin{corollary}[A changing-filtration stopping consumer]
\label{cor:v18-stopping}
Let $\psi(t,x)\in[0,1]$ be measurable and $L_\psi$-Lipschitz in $x$,
uniformly in $t$. For a bounded latent $f$ with $|f|\le1$, put
$M_t^n=\E_{P_n}[f(Z)\mid\mathcal F_t]$. Then
\[
 \left|\sup_{\tau\le T}\E_{P_n}\psi(\tau,M_\tau^n)
       -\sup_{\tau\le T}\E_P\psi(\tau,M_\tau)\right|
 \le\varepsilon_n+L_\psi\sqrt{80\varepsilon_n}.
\]
The same statement for the entropic value $U$ has the additional
Lipschitz factor $(e^{2\gamma}-1)/\gamma$ on the square-root term.
The suprema range over the common canonical stopping times; attainment
is not needed for the estimate.
\end{corollary}
\begin{proof}
For a fixed stopping time, first change the measure while keeping its
$[0,1]$ reward function, at cost at most $\varepsilon_n$. Then compare
the two reward functions on the same canonical path under $P$, using
Theorem~\ref{thm:v17-optional}, Cauchy--Schwarz and the Lipschitz bound.
The estimate is uniform in the stopping time, so it survives both
suprema. Equivalence of the observation laws is what allows the same
stopping-time class. The entropic assertion follows from the Lipschitz
calculation in the preceding theorem.
\end{proof}

\subsection{Resource transport is not the posterior coupling}
Apply any common finite causal acquisition map to $Y$ and retain a fixed
finite mark $w(Z)$. Data processing makes its marked row variations at
most $\varepsilon_n$, uniformly over a declared finite feedback family.
If the source and target have respective bounds $\alpha_n,\beta_n$,
Theorem~\ref{thm:v13-main} gives same-budget deficiency perturbation at
most $\alpha_n+\beta_n$, and
Proposition~\ref{prop:v18-audittransport} transports a fixed reset audit.
A subsequent consumer of width $S$ composed with a simulator of width $K$
has width $KS$, with the stated selector price. None of these executable
claims uses the latent-preserving maximal coupling. That coupling belongs
to the proof of conditional-process convergence, not to the finite-state
machine.

These results provide a direct changing-trajectory, likelihood,
Girsanov and bounded entropic-value input for a conditional-law theory.
They do not assert the strict-dual, rigidity, covariant-form or full kinetic
conclusions of a larger program without their separate analytic proofs.

\begin{example}[Why an amplitude bound is not a tightness argument]
\label{ex:v18-fiditightness}
Let $\xi_1,\xi_2$ be independent symmetric signs, and, for $n\ge5$, let
$\mathcal G_t^n$ reveal $\xi_1$ at $1/2-2/n$ and $\xi_2$ at $1/2-1/n$.
On $[0,1]$ the bounded martingale
\[
 M_t^n=\E[\xi_1+\xi_2\mid\mathcal G_t^n]
\]
converges at every finite collection of fixed times to the martingale
$M_t=\1_{[1/2,1]}(t)(\xi_1+\xi_2)$. Nevertheless its laws are not tight
in the Skorokhod $J_1$ topology. On the event $\xi_2=-\xi_1$, of probability
$1/2$, its path has an excursion of height one on an interval of length
$1/n$ and is zero outside that interval. The three-time oscillation modulus
of such a path is one at every scale greater than $2/n$. This violates the
necessary $J_1$ tightness modulus criterion. All maximal moments are
uniformly bounded by $2^p$.

This example concerns the inference from finite-dimensional convergence
and a maximal amplitude estimate alone. It does not claim that every
possible stronger conditional-kernel hypothesis fails. The uniform
estimates in Theorem~\ref{thm:v18-gaussian} remove the need for that
inference in the observation class treated here.
\end{example}
